% =========================================================================
% Capítulo 2: Cinética de la partícula
% Tema 01: Dinámica de la partícula
% Problema 18: Partícula con resorte en brazo giratorio
% =========================================================================
\subsubsection{Partícula con resorte en brazo giratorio}

Una partícula de masa $m$ está unida a un resorte lineal con constante elástica $K$ y longitud natural $r_0$. El otro extremo del resorte está conectado, en el punto $P$, al extremo libre de un brazo rígido sin masa de longitud $l$. El brazo está articulado en su otro extremo $O$ y gira en una trayectoria circular con una rapidez angular constante $\omega$. Sabiendo que la gravedad actúa hacia abajo y definiendo $r$ como la longitud del resorte y $\theta$ como el ángulo de la partícula respecto a la vertical, realice lo siguiente:

\begin{enumerate}
    \item Defina los marcos de referencia y sistemas de coordenadas adecuados.
    \item Analice la cinemática de la partícula determinando sus vectores de posición y velocidad en función de las coordenadas generalizadas.
    \item Analice la cinética del sistema para determinar las ecuaciones diferenciales del movimiento.
\end{enumerate}

\begin{figure}[!ht]
    \centering
    \begin{tikzpicture}[>=Latex, scale=1.5]
        % --- Parámetros ---
        \def\ang{-40}      % Ángulo de inclinación (omega*t y theta coinciden en el dibujo)
        \def\Lbar{2.0}     % Longitud del brazo l
        \def\Lspring{1.5}  % Longitud del resorte r

        % Coordenadas
        \coordinate (O) at (0,0);
        \coordinate (P) at (\ang:\Lbar);
        \coordinate (M) at (\ang+15:{\Lbar+\Lspring});

        % --- 1. Soporte y Pivote ---
        \fill[pattern=north east lines] (-0.5, 0) rectangle (0.5, 0.2);
        \draw[thick] (-0.5, 0) -- (0.5, 0);
        \filldraw[gray] (O) circle (2pt);
        \node[below left, xshift=-3pt] at (O) {$O$};

        % --- 2. Brazo Giratorio ---
        \draw[line width=3.5pt, black,line cap=round] (O) -- (P);
        \draw[line width=3pt, gray!40,line cap=round] (O) -- (P);
        \draw[thick,line cap=round ] (O) -- (P) node[midway, above right] {$l$};
        \filldraw (P) circle (0.5pt) node[below left, yshift=5pt] {$P$};
        \filldraw (O) circle (0.5pt);

        % --- 3. Resorte y Masa ---
        \draw[decoration={coil, aspect=0.4, segment length=3pt, amplitude=3pt, pre length=3pt, post length=3pt}, decorate, thick] (P) -- (M);
        \node at ($(P)!0.5!(M)$) [below left=2pt] {$K$};
        
        \shadedraw[ball color=white] (M) circle (3.5pt) node[] {$m$};

        % --- 4. Cotas y Referencias ---
        % Referencia Vertical
        \draw[dashed, gray] (O) -- (0, -2.0);
        \draw[dashed, gray] (P) -- ++ (0, -2.0);
        
        % Ángulo omega*t (Referido al brazo)
        \draw[->, thick] (0, -0.8) arc (-90:\ang:0.8);
        \node at ({(\ang-90)/2}:1.1) {$\omega t$};

        % Ángulo theta (Referido a la masa - dibujado un poco más abajo)
        \draw[->, thick] (P)++(0,-1.6) arc (-90:\ang+30:1.6) node[midway, above] {$\theta$};

        % Cota r (Longitud del resorte)
        % Calculamos puntos desplazados para la cota paralela
        \draw[thin, |<->|] ($(P)+(\ang+90:0.4)$) -- ($(M)+(\ang+90:0.4)$) node[midway, above right] {$r$};

    \end{tikzpicture}
    \caption{Sistema de partícula unida a un resorte acoplado a un brazo en rotación forzada.}
    \label{fig:resorte_brazo_giratorio}
\end{figure}